% AulaTeX - Actividad 9 construida desde la planeación, presentación y contenidos S9.
% Tema: Sistemas de ahorro para el retiro y retos de cobertura efectiva.
% Rúbrica: diagnóstico, fundamentos, análisis y propuestas (meta: 10/10).
\documentclass[11pt,letterpaper,oneside]{article}
\usepackage[margin=2.2cm]{geometry}
\usepackage[utf8]{inputenc}
\usepackage[T1]{fontenc}
\usepackage[spanish,es-tabla]{babel}
\usepackage[scaled=.96]{helvet}
\renewcommand{\familydefault}{\sfdefault}
\usepackage{setspace}
\usepackage{graphicx}
\usepackage[table]{xcolor}
\usepackage{float}
\usepackage{pdflscape}
\usepackage{tikz}
\usetikzlibrary{arrows.meta,positioning,fit,shapes.geometric,backgrounds}
\usepackage{hyperref}
\usepackage[authoryear,round]{natbib}
\usepackage[most]{tcolorbox}
\usepackage{attachfile}
\usepackage{enumitem}
\usepackage{listings}

% Estilo de bloque para reproducir la participación textual del foro
\lstdefinestyle{foroplano}{%
  basicstyle=\ttfamily\small\color{unadmInk},
  backgroundcolor=\color{unadmPaper},
  breaklines=true,
  columns=fullflexible,
  keepspaces=true,
  numbers=left,
  numberstyle=\tiny\color{unadmGreen},
  frame=leftline,
  framerule=1.4pt,
  rulecolor=\color{unadmGreen},
  xleftmargin=1.4em,
  extendedchars=true,
  inputencoding=utf8,
  literate=
    {á}{{\'a}}1 {é}{{\'e}}1 {í}{{\'i}}1 {ó}{{\'o}}1 {ú}{{\'u}}1
    {Á}{{\'A}}1 {É}{{\'E}}1 {Í}{{\'I}}1 {Ó}{{\'O}}1 {Ú}{{\'U}}1
    {ñ}{{\~n}}1 {Ñ}{{\~N}}1 {¿}{{?`}}1 {¡}{{!`}}1
}

% attachfile: el 'color' debe darse como TRIPLE RGB numérico (0..1)
\attachfilesetup{color={0.373 0.561 0.227},print=false}

% Botón de copia: appearance={} oculta el PushPin y muestra el 2.o argumento como
% botón visible.
\newcommand{\foroCopyButton}[1]{%
  \textattachfile[appearance={},description={Participacion del foro (texto copiable)},mimetype=text/plain]{#1}{%
    \colorbox{unadmGreen}{\small\textcolor{white}{\,Copiar participación\,}}%
  }%
}

\definecolor{unadmGreenDark}{HTML}{174A3A}
\definecolor{unadmGreen}{HTML}{5F8F3A}
\definecolor{unadmGold}{HTML}{B88A2A}
\definecolor{unadmPaper}{HTML}{F6F7F2}
\definecolor{unadmInk}{HTML}{1F2A24}
\definecolor{unadmBlue}{HTML}{DDEBE8}
\definecolor{unadmSand}{HTML}{F3E9D2}

% Entorno visual del foro (banda verde, fondo claro, texto justificado)
\newtcolorbox{forobox}[1][]{%
  enhanced, breakable,
  colback=unadmPaper, colframe=unadmGreen,
  boxrule=0pt, leftrule=3pt, arc=1.2mm,
  left=4mm, right=4mm, top=3mm, bottom=3mm,
  fonttitle=\bfseries\color{white}, coltitle=white, colbacktitle=unadmGreenDark,
  attach boxed title to top left={xshift=4mm,yshift=-2mm},
  boxed title style={colframe=unadmGreenDark,arc=0.6mm},
  #1
}
\hypersetup{colorlinks=true,linkcolor=unadmGreenDark,citecolor=unadmGreenDark,urlcolor=unadmGreenDark}
\setlength{\parindent}{0pt}
\setlength{\parskip}{0.48em}
\setstretch{1.12}
\setlength{\emergencystretch}{2em}
\raggedbottom
\bibliographystyle{plainnat}
\newcommand{\documenttitle}{Sistemas de ahorro para el retiro y retos de cobertura efectiva}
\newcommand{\documentsubtitle}{Diagnóstico y propuestas de mejora}
\newcommand{\documentsubject}{Actividad 9 - Derecho a la Seguridad Social}

\begin{document}

\begin{titlepage}
  \pagecolor{unadmPaper}\color{unadmInk}
  \begin{center}
    \vspace*{0.55cm}
    \includegraphics[height=1.95cm]{img/departamentos/UnADM.pdf}\\[0.55cm]
    {\Large Universidad Abierta y a Distancia de México}\\[0.18cm]
    {\large Licenciatura en Derecho}\\[1cm]
    {\color{unadmGreenDark}\rule{0.82\linewidth}{1.4pt}}\\[0.5cm]
    {\Huge\bfseries \documenttitle\par}
    \vspace{0.3cm}
    {\Large \documentsubtitle\par}
    \vspace{0.12cm}
    {\large \documentsubject\par}
    \vspace{0.4cm}
    {\color{unadmGold}\rule{0.62\linewidth}{1.1pt}}\\[0.85cm]
    \begin{tabular}{rl}
      \textbf{Alumno:} & Martin Jonathan de la Cruz \\
      \textbf{Matrícula:} & ES2611202040 \\
      \textbf{Figura docente:} & Cristian Ortega Barrera \\
      \textbf{Materia:} & Derecho a la Seguridad Social \\
      \textbf{Semestre/Bloque:} & 2 / 1 \\
      \textbf{Unidad temática:} & Sistemas de ahorro para el retiro \\
      \textbf{Fecha de entrega:} & 06 de septiembre de 2026
    \end{tabular}
    \vfill
    {\small Ciudad de México}
  \end{center}
  \nopagecolor
\end{titlepage}

\section{Introducción}

Este documento desarrolla la Actividad 9 con base en la planeación y los contenidos de la semana 9 (archivos locales del curso) \citep{materialS9DSS2026}. Su objetivo es ofrecer un diagnóstico claro sobre los retos de cobertura en los sistemas de ahorro para el retiro en México, comparar los efectos de los distintos regímenes y proponer medidas concretas para mejorar la inclusión y la suficiencia pensionaria.

La argumentación combina: (i) los materiales didácticos de la semana 9 y las actividades previas; (ii) la normativa aplicable (LSS, LISSSTE); y (iii) literatura analítica sobre el sistema de ahorro para el retiro \citep{lss2026,lissste2026,consarSitio2026,bricenoRuiz2011}.

\section{Desarrollo: retos de cobertura y propuestas}

En esta sección se presenta la intervención desarrollada en el foro sobre los temas 2 a 4 (diagnóstico, comparación de regímenes y propuestas). El título de la sección es breve y el desarrollo completo aparece dentro del bloque de foro que sigue, para facilitar la lectura y la copia de la participación.

\begin{forobox}[title={Participación desarrollada — Temas 2 a 4}]
\hfill\foroCopyButton{foro-participacion-Actividad-9.txt}\\[-0.5em]
\phantomsection\label{foro-participacion-9}

\textbf{Presentación}

Hola a todas y todos: comparto a continuación mi intervención para el foro sobre los sistemas de ahorro para el retiro. He estructurado la participación con un diagnóstico conciso, un análisis comparado y propuestas de política pública, y al final respondo las preguntas guía con sugerencias prácticas.

\medskip
\textbf{Diagnóstico ampliado}

La cobertura efectiva del retiro en México enfrenta tres problemas integrados: (i) la persistente informalidad laboral y el subregistro, que impiden la acumulación continua de cotizaciones; (ii) la insuficiencia de aportaciones y los costos administrativos (comisiones, bajos rendimientos reales) que erosionan los saldos acumulados; y (iii) la fragmentación institucional entre regímenes que dificulta la portabilidad y la coordinación de políticas. Estos elementos se combinan para producir bajos niveles de suficiencia pensionaria, especialmente entre mujeres y personas con carreras discontinuas por responsabilidades de cuidado \citep{consarSitio2026,bricenoRuiz2011}.

\medskip
\textbf{Análisis comparado breve}

Las cuentas individuales mejoran la transparencia y la trazabilidad de aportes, pero dependen críticamente de la continuidad contributiva y de costos operativos bajos para ser eficaces. Los regímenes públicos y los mecanismos redistributivos (complementos, pensión mínima) son herramientas necesarias para corregir insuficiencias que la lógica contributiva no resuelve por sí sola. En la práctica, la combinación de instrumentos (microcuentas, complementos y políticas laborales) muestra mejores resultados de inclusión que la aplicación aislada de cualquiera de ellos.

\medskip
\textbf{Propuestas concretas (implementación práctica)}
\begin{enumerate}
  \item \textbf{Microcuentas universales con bonificación temporal:} habilitar cuentas de baja fricción para trabajadores informales con aportaciones mínimas y una bonificación estatal en los primeros 24 meses para incentivar la adhesión.
  \item \textbf{Reducción y transparencia de costos:} establecer límites regulatorios a comisiones y publicar comparadores públicos de rendimiento neto; promover subastas o modelos colectivos de inversión para reducir gastos administrativos.
  \item \textbf{Complementos focalizados:} garantizar una pensión mínima universal o focalizada mediante complementos que se activen cuando los saldos acumulados sean insuficientes, evitando incentivos adversos si van acompañados de requisitos de contratación mínima.
  \item \textbf{Protección para carreras discontinuas:} reconocer periodos de cuidado para efectos de cuentas, ofrecer cotizaciones bonificadas y facilitar aportaciones voluntarias con co-financiamiento público para mujeres y cuidadores.
  \item \textbf{Registro y formalización digital:} integrar registros fiscales y de seguridad social para facilitar altas y bajas automáticas y detectar subregistro; acompañarlo con campañas de capacitación y alivios fiscales temporales.
\end{enumerate}
% La respuesta a una participación se reubicó en un forobox independiente y tiene su propio archivo .txt

\end{forobox}

\begin{forobox}[title={Respuesta a una participación}]
\hfill\foroCopyButton{foro-respuesta-Actividad-9.txt}\\[-0.5em]
\textbf{Respuesta}

Gracias por tu aporte; comparto y complemento brevemente mis respuestas a tus planteamientos.

Primero, coincido en que la barrera central para que las cuentas individuales generen pensiones adecuadas es la continuidad de las cotizaciones. La informalidad, la movilidad laboral y los periodos sin aportaciones fragmentan la historia contributiva y reducen significativamente la probabilidad de alcanzar una pensión suficiente. Además, las comisiones y los bajos rendimientos reales actúan como factores que erosionan los saldos acumulados y agravan la insuficiencia.

Segundo, respecto a medidas prácticas para aumentar la cobertura, tu propuesta sobre microcuentas me parece adecuada; añadiría tres precisiones operativas:
\begin{itemize}
  \item Empezar con una fase piloto regional y un periodo de bonificación estatal (18–24 meses) para generar evidencia y ajustar el diseño.
  \item Combinar la bonificación con campañas de registro digital que reduzcan la fricción administrativa y faciliten la identificación de beneficiarios potenciales.
  \item Asegurar que las microcuentas vayan acompañadas de medidas para reducir comisiones (por ejemplo, modelos colectivos de inversión) y de un complemento público que actúe como red de seguridad cuando los saldos no alcancen la suficiencia.
\end{itemize}

En resumen: coincido en el diagnóstico y propongo articular microcuentas, formalización digital y límites a comisiones como un paquete coordinado; si te parece útil, puedo preparar un esquema operativo breve con indicadores de seguimiento.

Quedo atento a tus comentarios.
\end{forobox}

\section{Conclusiones}

La intervención desarrollada y el análisis completo aparecen en la Sección~\ref{foro-participacion-9} (bloque de foro) y en el archivo adjunto `foro-participacion-Actividad-9.txt`. En síntesis, la eficacia del sistema requiere acciones combinadas: formalización laboral, reducción de costos administrativos, complementos redistributivos y medidas específicas para carreras discontinuas.

\textbf{Declaración de uso de inteligencia artificial.} Se utilizó asistencia de IA para organizar y sintetizar el contenido; el autor revisó y validó las referencias y asume la responsabilidad final sobre el texto.

\nocite{bricenoRuiz2011,lissste2026,lss2026,consarSitio2026,materialS9DSS2026}
\bibliography{derecho-a-la-seguridad-social}

\end{document}
